
%%%%%%%%%%%%%%%%%%%%%%%%%%%%%%%%%%%%%%%%%
% University Assignment Title Page 
% LaTeX Template
% Version 1.0 (27/12/12)
%
% This template has been downloaded from:
% http://www.LaTeXTemplates.com
%
% Original author:
% WikiBooks (http://en.wikibooks.org/wiki/LaTeX/Title_Creation)
%
% License:
% CC BY-NC-SA 3.0 (http://creativecommons.org/licenses/by-nc-sa/3.0/)
% 
% Instructions for using this template:
% This title page is capable of being compiled as is. This is not useful for 
% including it in another document. To do this, you have two options: 
%
% 1) Copy/paste everything between \begin{document} and \end{document} 
% starting at \begin{titlepage} and paste this into another LaTeX file where you 
% want your title page.
% OR
% 2) Remove everything outside the \begin{titlepage} and \end{titlepage} and 
% move this file to the same directory as the LaTeX file you wish to add it to. 
% Then add \input{./title_page_1.tex} to your LaTeX file where you want your
% title page.
%
%%%%%%%%%%%%%%%%%%%%%%%%%%%%%%%%%%%%%%%%%
%\title{Title page with logo}
%----------------------------------------------------------------------------------------
%	PACKAGES AND OTHER DOCUMENT CONFIGURATIONS
%----------------------------------------------------------------------------------------

\documentclass[12pt]{article}
\usepackage[spanish]{babel}
\usepackage[utf8x]{inputenc}
\usepackage{amsmath}
\usepackage{graphicx}
\usepackage[colorinlistoftodos]{todonotes}
\usepackage{float} %Este paquete es para que se usen los entornos de las tablas, de modo que se puedan fijar en alguna posición en particular.
\usepackage{hyperref} %Este paquete es para poder insertar hipervínculos en el documento, como por ejemplo direcciones de e-mail.
\usepackage{textcomp} %Paquete para poner los logotipos de registrado,copyright, etc.
\usepackage{listings} %Paquete para los códigos en cualquier lenguage.
\usepackage{amsfonts}
\usepackage{amssymb}
\usepackage{eufrak}
                      
\begin{document}

\begin{titlepage}

  \newcommand{\HRule}{\rule{\linewidth}{0.5mm}} % Defines a new command for the horizontal lines, change thickness here

  \center % Center everything on the page
 
  % ----------------------------------------------------------------------------------------
  % HEADING SECTIONS
  % ----------------------------------------------------------------------------------------

  \textsc{\LARGE Universidad Nacional del
    Litoral}\\[1.5cm] % Name of your university/college
  \textsc{\Large Facultad de Ingeniería y Ciencias Hídricas
  }\\[0.5cm] % Major heading such as course name
  \textsc{\large Proyecto Final de Carrera:
    ``Análisis e implementación de un método para la proyección de
    soluciones en problemas de interacción fluido-estructura con mallas no coincidentes''}\\[0.5cm] % Minor heading such as course title

  % ----------------------------------------------------------------------------------------
  % TITLE SECTION
  % ----------------------------------------------------------------------------------------

  \HRule \\[0.4cm]
  { \huge \bfseries Entregable N°1:
    Revisión bibliográfica básica}\\[0.4cm] % Title of your document
  \HRule \\[1.0cm]
 
  % ----------------------------------------------------------------------------------------
  % AUTHOR SECTION
  % ----------------------------------------------------------------------------------------

  \begin{minipage}{0.4\textwidth}
    \begin{flushleft} \large
      \emph{Autores:}\\
        Fabrizio J. \textsc{Piva} \\ % Your name
        Pablo S. \textsc{Vera} \\ % Your name
    \end{flushleft}
  \end{minipage}
  ~
  \begin{minipage}{0.4\textwidth}
    \begin{flushright} \large
      \emph{Director:} \\
      Dr. Gustavo \textsc{Ríos} % Supervisor's Name
      \emph{co-Director:} \\
      Dr. Luciano \textsc{Garelli} % Supervisor's Name
    \end{flushright}
  \end{minipage}\\[1cm]

  % If you don't want a supervisor, uncomment the two lines below and
  % remove the section above
  % \Large \emph{Autores:}\\
  % Fabrizio J. \textsc{Piva} \\ % Your name
  % Pablo S. \textsc{Vera} \\ [1.0cm]% Your name

  % ----------------------------------------------------------------------------------------
  % DATE SECTION
  % ----------------------------------------------------------------------------------------

  {\large
    Ingeniería Informática}\\[1cm] % Date, change the \today to a set date if you want to be precise
  
  % ----------------------------------------------------------------------------------------
  % LOGO SECTION
  % ----------------------------------------------------------------------------------------
  \begin{center}
    \includegraphics[width=0.33\textwidth]{logo_unl.eps}\\[1cm] % Include a department/university logo - this will require the graphicx package
  \end{center}
  % ----------------------------------------------------------------------------------------

  \vfill % Fill the rest of the page with whitespace

\end{titlepage}

\begin{abstract}
  El presente documento es un breve informe que contendrá, entre
  otros aspectos, una breve descripción de los conceptos que se
  adquirieron en la etapa del cronograma denominada ``Revision bibliografica básica''.
\end{abstract}

\section{Introducción}

Este informe tiene como objetivo desarrollar brevemente cuales han
sido los conocimientos adquiridos a lo largo de esta etapa, entre los
cuales se pueden mencionar: esquemas de discretización, métodos
numéricos, interacción fluído-estructura, conceptos de C++, entre
otros. A lo largo del mismo, se abordarán las diferentes temáticas
para luego finalizar con la muestra de algunos casos de prueba que
permitan demostrar la correcta integración de todos estos temas que se
abordaron a lo largo de esta etapa.

\section{Esquemas de discretización}
\label{sec:discretizacion}
%-------------------- ACA QUEDÉ !!!!! --------------------
Para realizar el algoritmo, basicamente debemos recordar que el
producto de dos matrices $A_{mxn}$ y $B_{nxp}$ dan como resultado una matriz
$C_{mxp}$, donde:

\begin{equation}
  c_{i,j} = \sum_{k=0}^n a_{i,k} \cdot b_{k,j} 
\end{equation}

Teniendo en cuenta lo anteriormente mencionado, y considerando que se
utilizó la librería Eigen para facilitar la indexación y llenado de
matrices, en C++ puede traducirse como:

\lstset{language=C++, breaklines=true, basicstyle=\footnotesize}
\begin{lstlisting}[frame=single]
  #include <iostream>
  #include <Eigen/Dense>
  #include <time.h>

  using namespace Eigen;
  using namespace std;

  MatrixXf  multiplicar_matrices( const MatrixXf &A,const MatrixXf &B ){

  int nfilas = A.rows();
  int ncols = B.cols();
  int limk = A.cols();
  MatrixXf C = MatrixXf::Zero(nfilas, ncols );
  
  for (int i=0;i<nfilas ; ++i){
    for (int j=0; j<ncols ; ++j) {
      for (int k=0; k<limk ; ++k) {
	C(i,j) += A(i,k)*B(k,j);
      }  
    }
  }
 return C;
}

\end{lstlisting}

Consideremos que para todo este ejercicio, utilizaremos como prueba
dos matrices de 3x3. Entonces cuando ejecutamos este algoritmo,  el
tiempo de ejecución es de $1.2\cdot10^{-5}$ segundos. En las
siguientes subsecciones veremos si este resultado mejora alterando el
orden de los lazos. 

\subsection{Modificación 1: lazo interno recorre las filas}
\label{sec:matricesv1}

Esta modificación consiste básicamente en cambiar el algoritmo
estándar de multiplicación de matrices para que el lazo interno sea
aquel que recorre el índice de filas. Esto en C++ puede ponerse en
práctica simplemente alterando el orden de los ciclos \emph{for}, como
se muestra en la función \emph{multiplicacion\_matrices\_v1}:

\lstset{language=C++, breaklines=true, basicstyle=\footnotesize}
\begin{lstlisting}[frame=single]

  MatrixXf  multiplicar_matrices_v1( const MatrixXf &A,const MatrixXf &B ){

  int nfilas = A.rows();
  int ncols = B.cols();
  int limk = A.cols();
  MatrixXf C = MatrixXf::Zero(nfilas, ncols );
 
  for (int j=0; j<ncols ; ++j) {
    for (int k=0; k<limk ; ++k) {
      for (int i=0; i<nfilas ; ++i) {
        C(i,j) += A(i,k)*B(k,j);
      }  
    }
  }
  return C;
  }

\end{lstlisting}

Prestar especialmente atención, que si comparamos el algoritmo
estándar con esta \emph{version 1}, vemos claramente que en esta nueva
versión el lazo interno recorre las filas. Esta modificación produce
un impacto significativo en el tiempo de ejecución, ya que para esta
nueva versión es de $9\cdot10^{-06}$ segundos.

\subsection{Modificación 2: lazo interno recorre las columnas}
\label{sec:matricesv2}

Esta última modificación consiste básicamente en cambiar el algoritmo
estándar de multiplicación de matrices para que el lazo interno sea
aquel que recorre el índice de columnas. Como ya se mencionó en la
sección \ref{sec:matricesv1}, alterando los ciclos \emph{for} se
obtiene la función \emph{multiplicacion\_matrices\_v2}:

\lstset{language=C++, breaklines=true, basicstyle=\footnotesize}
\begin{lstlisting}[frame=single]

MatrixXf  multiplicar_matrices_v2( const MatrixXf &A,const MatrixXf &B){
								 
  int nfilas = A.rows();
  int ncols = B.cols();
  int limk = A.cols();
  MatrixXf C = MatrixXf::Zero(nfilas, ncols );
 
  for (int i=0; i<nfilas ; ++i) {
    for (int k=0; k<limk ; ++k) {
      for (int j=0; j<ncols ; ++j) {
	C(i,j) += A(i,k)*B(k,j);
      }  
    }
  }
 return C;
}

\end{lstlisting}

Asimismo, esta versión que coloca el lazo interno que recorra las
columnas, produce un impacto significativo en el tiempo de ejecución
si lo comparamos con el del algoritmo de multiplicación estándar, ya que
para esta nueva versión es de $9\cdot10^{-06}$ segundos.

\section{Aplicación de desenrollado de ciclos}
\label{sec:desenrollado}

En esta sección se pondrá en manifiesto como aplicar la técnica de
desenrollado de ciclos en un algoritmo. La misma consiste en, dado un
ciclo iterativo común y corriente, disminuir la cantidad de
repeticiones mediante la codificación explícita de algunas o todas las
iteraciones. Esta técnica produce una mejora significativa en el
tiempo de ejecución del algoritmo, como lo veremos en el siguiente
ejemplo.

Supongamos que tenemos una matriz de Nx3, donde cada fila contiene un
vector en el espacio tridimensional. Ahora supongamos que necesitamos
un algoritmo que calcule el módulo de cada vector que está almacenado
en esta matriz, y que el resultado se almacene en un vector de
Nx1. Este algoritmo podría implementarse en C++ de la siguiente
manera:

\lstset{language=C++, breaklines=true, basicstyle=\footnotesize}
\begin{lstlisting}[frame=single]

#include <iostream>
#include <Eigen/Dense>
#include <time.h>

using namespace Eigen;
using namespace std;

VectorXf modulos_vectores (MatrixXf M){
  int nfilas = M.rows();
  VectorXf R = VectorXf::Zero(nfilas,1);
  
  for (int i=0; i< nfilas; ++i){
    for (int j=0; j<3 ; ++j) {
      R(i,0) += M(i,j)*M(i,j);
    }
    R(i,0) = sqrt(R(i,0) );
  }
  return R;
}

\end{lstlisting}
 
Para aplicar el método de desenrollado de ciclos, podríamos
deshacernos del ciclo interno simplemente escribiendolo de manera
explícita como se hizo en la función
\emph{modulos\_vectores\_desenrollado}:

\lstset{language=C++, breaklines=true, basicstyle=\footnotesize}
\begin{lstlisting}[frame=single]

VectorXf modulos_vectores_desenrollado (MatrixXf M){
  int nfilas = M.rows();
  VectorXf R = VectorXf::Zero(nfilas,1);
  
  for (int i=0; i< nfilas; ++i){
      R(i,0) = M(i,0)*M(i,0) + M(i,1)*M(i,1) + M(i,2)*M(i,2);
      R(i,0) = sqrt( R(i,0) );
    }

  return R;
}

\end{lstlisting}

A continuación se hará una prueba comparativa de estos dos algoritmos,
con matrices de tamaño 10, 100, y 1000. Los resultados de los tiempos
de ejecución pueden verse en la tabla \ref{tab:resultados}, donde
podemos observar que el desenrollado tiene un desempeño
notablemente mejor considerando el tiempo de ejecución.

\begin{table}[H]
  \centering
  \begin{tabular}{l|r}
    Sin desenrollado & Con desenrollado \\\hline
    1.9E-05 & 1E-05 \\
    1.13E-04 & 8.1E-05 \\
    8.71E-04 & 6.57E-04
  \end{tabular}
  \caption{\label{tab:resultados}Tabla comparativa de tiempos de
    ejecución en segundos.}
\end{table}

\section{Conclusiones}

Haciendo un análisis de los resultados obtenidos en las secciones
\ref{sec:matrices} y \ref{sec:desenrollado} podemos concluir que las
técnicas de reordenamiento y desenrollado de ciclos son muy útiles a
la hora de programar, ya que en este informe se ha puesto en evidencia
mediante sencillos ejemplos que los tiempos de ejecución mejoran
notablemente cuando se utilizan las mismas. Además, estos tiempos de
ejecución podrían mejorarse aun mas si se aplicaran ambas técnicas de
manera simultánea o si se complementaran con otras técnicas que se han
visto a lo largo de esta unidad temática.

 \end{document}


% \begin{figure}
 %   \centering
%   \includegraphics[width=1.0\textwidth]{resultados.eps}
%   \caption{\label{fig:reporte-ORNL}Reporte a enviar a la ORNL: ejemplo
%   1000d}
% \end{figure}

% \subsection{Tables and Figures}

% Use the table and tabular commands for basic tables --- see
% Table~\ref{tab:widgets}, for example. You can upload a figure (JPEG,
% PNG or PDF) using the files menu. To include it in your document,
% use the includegraphics command as in the code for
% Figure~\ref{fig:frog} below.

% % % Commands to include a figure:
% % \begin{figure}
% % \centering
% % \includegraphics[width=0.5\textwidth,natwidth=399,natheight=356]{curva.jpg}
% % \caption{\label{fig:frog}This is a figure caption.}
% % \end{figure}

% % \begin{table}
% % \centering
% % \begin{tabular}{l|r}
% % Item & Quantity \\\hline
% % Widgets & 42 \\
% % Gadgets & 13
% % \end{tabular}
% % \caption{\label{tab:widgets}An example table.}
% % \end{table}

% \subsection{Matematicas}

% \LaTeX{} is great at typesetting mathematics. Let $X_1, X_2, \ldots,
% X_n$ be a sequence of independent and identically distributed random
% variables with $\text{E}[X_i] = \mu$ and $\text{Var}[X_i] = \sigma^2
% < \infty$, and let
% $$S_n = \frac{X_1 + X_2 + \cdots + X_n}{n}
% = \frac{1}{n}\sum_{i}^{n} X_i$$ denote their mean. Then as $n$
% approaches infinity, the random variables $\sqrt{n}(S_n - \mu)$
% converge in distribution to a normal $\mathcal{N}(0, \sigma^2)$.

% \subsection{Lists}

% You can make lists with automatic numbering \dots

% \begin{enumerate}
% \item Like this,
% \item and like this.
% \end{enumerate}
% \dots or bullet points \dots
% \begin{itemize}
% \item Like this,
% \item and like this.
% \end{itemize}

% We hope you find write\LaTeX\ useful, and please let us know if you
% have any feedback using the help menu above.

